% !TEX program = xelatex
\documentclass[12pt,a4paper]{article}

\usepackage{ctex}
\usepackage[utf8]{inputenc}
\usepackage{fontspec}
\setmainfont{Times New Roman}

\usepackage[top=2.5cm, bottom=2.5cm, left=3cm, right=3cm]{geometry}
\usepackage{setspace}
\onehalfspacing

\usepackage{amsmath, amssymb}
\usepackage{graphicx}
\usepackage{hyperref}
\usepackage{booktabs}
\usepackage{array}
\usepackage{listings}
\usepackage{xcolor}
\usepackage{fancyhdr}
\usepackage{titlesec}
\usepackage{caption}
\usepackage{longtable}
\usepackage{float}

\lstset{
    basicstyle=\small\ttfamily,
    backgroundcolor=\color{gray!10},
    frame=single,
    numbers=left,
    numberstyle=\tiny,
    breaklines=true,
    showstringspaces=false,
    language=Python,
    keywordstyle=\color{blue},
    commentstyle=\color{green!50!black},
    stringstyle=\color{red}
}

\pagestyle{fancy}
\fancyhf{}
\fancyhead[C]{遥感图像分析平台 -- 项目主报告}
\fancyfoot[C]{\thepage}

\titleformat{\section}{\Large\bfseries}{\thesection}{1em}{}
\titleformat{\subsection}{\large\bfseries}{\thesubsection}{1em}{}
\titleformat{\subsubsection}{\normalsize\bfseries}{\thesubsubsection}{1em}{}

\begin{document}

% ── 封面 ──
\begin{titlepage}
\centering
\vspace*{3cm}
{\Huge\bfseries 遥感图像分析平台}\\[0.8cm]
{\large 基于轻量级深度学习的建筑分割与变化检测系统}\\[1.5cm]
{\large \textbf{《人工智能基础B》课程项目主报告}}\\[0.5cm]
\vfill
\begin{tabular}{rl}
\textbf{小组成员：} & Air-0000、Member2、Member3 \\
\textbf{完成日期：} & 2026年6月 \\
\end{tabular}
\vspace{2cm}
\end{titlepage}

\newpage
\tableofcontents
\newpage

%====================================================================
\section{项目概述（占比10\%）}
%====================================================================

\subsection{项目背景}

遥感技术通过卫星或无人机平台获取地表信息，是城市规划、环境监测、灾害评估等领域的关键数据来源。近年来，高分辨率遥感卫星（如 WorldView、GeoEye）和无人机航拍技术的普及使得遥感图像数据量呈指数级增长，如何从海量遥感数据中自动提取有用信息成为亟待解决的问题。

\textbf{现有方案痛点分析：}

\begin{enumerate}
    \item \textbf{模型参数量过大：} 主流分割模型（DeepLabV3+、PSPNet）参数量在 40M--100M，训练和推理需要高性能 GPU，普通用户难以复现；
    \item \textbf{依赖预训练权重：} 多数方法依赖 ImageNet 预训练，增加了额外存储和领域迁移成本，且自然图像与遥感图像存在域差异；
    \item \textbf{部署门槛高：} 复杂的环境配置和依赖管理阻碍了非专业用户的使用；
    \item \textbf{缺少一站式工具：} 现有方案多聚焦单一任务（分割或检测），缺乏统一平台。
\end{enumerate}

\subsection{项目目标}

基于 Python + PyTorch 深度学习框架，构建一个轻量级遥感图像分析平台，实现建筑语义分割与双时相变化检测两大功能。\textbf{可量化考核指标：}

\begin{itemize}
    \item 语义分割 mIoU（平均交并比）$\geq$ 85\%；
    \item 变化检测 F1 分数 $\geq$ 60\%；
    \item 单图推理耗时 $\leq$ 50ms（GPU）；
    \item 端到端 Web 交互界面，支持图像上传与实时结果展示。
\end{itemize}

\subsection{团队分工}

\begin{table}[H]
\centering
\caption{团队分工}
\label{tab:team}
\begin{tabular}{cll}
\toprule
成员 & 职责 & 具体工作内容 \\
\midrule
Air-0000 & 方案设计与算法实现 & 模型架构设计、训练脚本编写、超参数调优 \\
Member2 & 数据处理与界面开发 & 数据集加载模块、Streamlit 交互界面、可视化 \\
Member3 & 报告撰写与测试验证 & 实验记录、性能测试、报告与视频录制 \\
\bottomrule
\end{tabular}
\end{table}

%====================================================================
\section{总体方案设计（占比20\%）}
%====================================================================

\subsection{总体架构图}

系统自底向上分为数据层、模型层和交互层：

\begin{lstlisting}[frame=none,numbers=none,language={}]
交互层
  Streamlit Web 界面
    - 语义分割模式 (单图上传 + 实时推理 + 面积统计)
    - 变化检测模式 (双时相上传 + 概率图 + 统计)
        |
模型层
    - Tiny U-Net (2M)   建筑语义分割    输入: 256x256x3
    - FC-Siam-diff (625K) 双时相变化检测  输入: T1+T2 256x256x3
        |
数据层
    - WHU Building (分割) 4736/1036 训练/验证
    - LEVIR-CD (变化检测) 445/64/128 训练/验证/测试
\end{lstlisting}

完整业务流程图：

\begin{lstlisting}[frame=none,numbers=none,language={}]
原始图像 -> 统一尺寸 -> 归一化 -> 数据增强 -> DataLoader
       -> 模型训练 -> 评估 -> 权重保存
       -> 模型加载 -> Streamlit 推理 -> 结果可视化
\end{lstlisting}

\subsection{技术选型说明}

\begin{table}[H]
\centering
\caption{技术选型理由}
\label{tab:tech}
\begin{tabular}{lll}
\toprule
技术 & 用途 & 选型理由 \\
\midrule
Python 3.12 & 开发语言 & 生态丰富，AI/ML 库支持完善 \\
PyTorch 2.12+cu130 & 深度学习框架 & 动态图机制方便调试，社区活跃 \\
Streamlit & Web 界面 & 纯 Python 构建，无需前后端分离开发 \\
Albumentations & 数据增强 & 比 torchvision 快 3--5$\times$ \\
OpenCV & 图像处理 & 性能优异，支持多种格式 \\
CUDA 13.3 & GPU 加速 & 适配 RTX 5060 Ti (sm\_120) \\
\bottomrule
\end{tabular}
\end{table}

\subsection{模块划分}

\begin{table}[H]
\centering
\caption{功能模块划分}
\label{tab:modules}
\begin{tabular}{clll}
\toprule
模块 & 功能 & 技术栈 & 负责人 \\
\midrule
数据加载模块 & LEVIR-CD / WHU 数据读取 & Python, OpenCV & Member2 \\
数据增强模块 & 实时增强（翻转/旋转/归一化） & Albumentations & Air-0000 \\
模型定义模块 & Tiny U-Net / FC-Siam-diff & PyTorch & Air-0000 \\
训练模块 & 训练/验证循环、指标记录 & PyTorch, TensorBoard & Air-0000 \\
推理模块 & 模型加载、前向传播、后处理 & PyTorch, NumPy & Air-0000 \\
界面模块 & 文件上传、结果展示、交互控制 & Streamlit & Member2 \\
测试模块 & 功能/性能/稳定性测试 & Python, time & Member3 \\
\bottomrule
\end{tabular}
\end{table}

%====================================================================
\section{详细实现过程（占比40\%）}
%====================================================================

\subsection{数据模块}

\textbf{数据集来源：}

本项目使用两个公开遥感数据集：

\begin{itemize}
    \item \textbf{WHU Building Dataset：} 武汉大学发布，包含 8,189 张 512$\times$512 卫星图像，覆盖新西兰 Christchurch 地区，空间分辨率 0.3--0.5 m/pixel，标注为建筑/非建筑二值掩码。图像内容包含高密度市中心到低密度郊区的各类建筑形态。
    \item \textbf{LEVIR-CD Dataset：} 武汉大学 LEVIR 实验室发布，包含 637 对 1024$\times$1024 卫星图像，覆盖美国德州 20 个城市，变化类型包括新建建筑、拆除建筑、道路建设等。
\end{itemize}

\textbf{数据预处理方法：}

\begin{enumerate}
    \item \textbf{统一尺寸：} 所有图像 Resize 到 256$\times$256（双线性插值），标注使用最近邻插值防止类别混淆；
    \item \textbf{归一化：} 使用 ImageNet 统计量（mean=[0.485,0.456,0.406]，std=[0.229,0.224,0.225]），与主流方法保持一致的输入分布；
    \item \textbf{格式兼容：} 分割数据加载器支持多目录名（images/image/A，masks/mask/label 等），适应不同数据集的组织方式。
\end{enumerate}

\textbf{数据增强：}

训练阶段使用 Albumentations 库进行实机数据增强：

\begin{lstlisting}
A.Compose([
    A.RandomRotate90(p=0.5),
    A.HorizontalFlip(p=0.5),
    A.VerticalFlip(p=0.5),
    A.Normalize(mean, std),
    ToTensorV2(),
])
\end{lstlisting}

验证和测试阶段仅做归一化，不做随机增强。

\textbf{数据集划分：}

\begin{table}[H]
\centering
\caption{数据集划分}
\label{tab:split}
\begin{tabular}{lcccc}
\toprule
数据集 & 训练集 & 验证集 & 测试集 & 图像尺寸 \\
\midrule
WHU Building & 4,736 & 1,036 & 2,417 & 512$\times$512 \\
LEVIR-CD & 445对 & 64对 & 128对 & 256$\times$256 \\
\bottomrule
\end{tabular}
\end{table}

\subsection{模型模块}

\textbf{Tiny U-Net（建筑语义分割）：}

U-Net 采用对称的编码器-解码器结构，编码器逐层下采样提取特征，解码器通过 Skip Connection 融合编码器各层特征恢复空间分辨率。Tiny 版本将编码器通道压缩为 16$\to$32$\to$64$\to$128，参数量约 2M。

核心代码：

\begin{lstlisting}
class TinyUNet(nn.Module):
    def __init__(self, in_channels=3, num_classes=1, base_ch=16):
        super().__init__()
        c = base_ch
        # 编码器 (4层)
        self.enc1 = DoubleConv(in_channels, c)
        self.enc2 = DoubleConv(c, c*2)
        self.enc3 = DoubleConv(c*2, c*4)
        self.enc4 = DoubleConv(c*4, c*8)
        self.pool = nn.MaxPool2d(2)
        # 瓶颈
        self.bottleneck = DoubleConv(c*8, c*16)
        # 解码器 (含Skip Connection)
        self.dec4 = DoubleConv(c*16+c*8, c*8)
        self.dec3 = DoubleConv(c*8+c*4, c*4)
        self.dec2 = DoubleConv(c*4+c*2, c*2)
        self.dec1 = DoubleConv(c*2+c, c)
        self.out = nn.Conv2d(c, num_classes, 1)

    def forward(self, x):
        e1 = self.enc1(x)
        e2 = self.enc2(self.pool(e1))
        e3 = self.enc3(self.pool(e2))
        e4 = self.enc4(self.pool(e3))
        b = self.bottleneck(self.pool(e4))
        d4 = self.dec4(torch.cat([self.up(b), e4], 1))
        d3 = self.dec3(torch.cat([self.up(d4), e3], 1))
        d2 = self.dec2(torch.cat([self.up(d2), e2], 1))
        d1 = self.dec1(torch.cat([self.up(d2), e1], 1))
        return self.out(d1)  # logits
\end{lstlisting}

\textbf{FC-Siam-diff（变化检测）：}

采用孪生网络架构，T1 和 T2 图像通过\textbf{共享权重}的编码器提取特征，在特征空间做差（L1 距离）后解码重建变化图。参数量 625K。

核心代码：

\begin{lstlisting}
class SiamDiffChangeDetector(nn.Module):
    def forward(self, img_A, img_B):
        def encode(x):
            f1 = self.e1(x)
            f2 = self.e2(self.p1(f1))
            f3 = self.e3(self.p2(f2))
            return f1, f2, f3
        a1, a2, a3 = encode(img_A)
        b1, b2, b3 = encode(img_B)
        # 特征差 (L1) + 解码
        x = self.d2(torch.cat([
            self.up(torch.abs(a3 - b3)),
            torch.abs(a2 - b2)], 1))
        x = self.d1(torch.cat([
            self.up(x), torch.abs(a1 - b1)], 1))
        return self.out(x)  # logits
\end{lstlisting}

\textbf{训练参数设置：}

\begin{table}[H]
\centering
\caption{训练超参数}
\label{tab:hp}
\begin{tabular}{lcc}
\toprule
参数 & 语义分割 & 变化检测 \\
\midrule
优化器 & AdamW & AdamW \\
学习率 & $1\times10^{-3}$ & $1\times10^{-3}$ \\
Weight Decay & $1\times10^{-4}$ & $1\times10^{-4}$ \\
Batch Size & 64 & 32 \\
Epochs & 50 & 100 \\
损失函数 & BCEWithLogitsLoss & BCEWithLogitsLoss \\
梯度裁剪 & max\_norm=5.0 & max\_norm=5.0 \\
学习率调度 & CosineAnnealingLR & CosineAnnealingLR \\
\bottomrule
\end{tabular}
\end{table}

\textbf{训练过程：}

训练循环包含以下关键步骤：

\begin{enumerate}
    \item 每个 epoch 遍历训练集，前向传播计算损失，反向传播更新权重；
    \item 梯度裁剪防止梯度爆炸；
    \item 每个 epoch 结束后在验证集上评估性能；
    \item 验证指标提升时保存最佳模型权重；
    \item 余弦退火调度器逐步降低学习率。
\end{enumerate}

\textbf{性能指标：}

\begin{table}[H]
\centering
\caption{模型性能指标}
\label{tab:perf}
\begin{tabular}{lcc}
\toprule
指标 & 语义分割 & 变化检测 \\
\midrule
准确率 (Accuracy) & 97.3\% & 95.1\% \\
精确率 (Precision) & 0.891 & 0.623 \\
召回率 (Recall) & 0.868 & 0.682 \\
F1 分数 & 0.879 & 0.651 \\
IoU / mIoU & \textbf{0.877} & \textbf{0.477} \\
\bottomrule
\end{tabular}
\end{table}

\subsection{大模型模块}

本项目以轻量级 CNN 模型为核心，未直接集成 LLM。但大模型可在以下场景中辅助遥感分析：

\begin{itemize}
    \item \textbf{结果解读：} 将分割/变化检测结果以自然语言描述输出（"检测到 15 栋新建建筑，主要集中在东部区域"）；
    \item \textbf{智能问答：} 用户可以通过自然语言查询分析结果（"哪个区域变化最大？"）；
    \item \textbf{报告生成：} 自动生成分析报告，包含统计数据和文字总结。
\end{itemize}

\textbf{提示词设计思路（以 GLM-4 为例）：}

\begin{lstlisting}
prompt = f"""
你是一个遥感图像分析助理。
根据以下变化检测结果，生成简要分析报告：
- 变化面积占比：{change_ratio:.1f}%
- 最大变化区域坐标：({max_x}, {max_y})
- 置信度分布：平均 {avg_conf:.2f}
请用一句话总结变化情况。
"""
\end{lstlisting}

\subsection{界面/交互模块}

使用 Streamlit 构建 Web 交互界面，支持两种工作模式：

\textbf{语义分割模式：}
\begin{itemize}
    \item 上传单张遥感图像或选择样例图像；
    \item 自动运行模型推理，显示原图与分割覆盖图对比；
    \item 实时统计建筑面积占比；
    \item 阈值滑块（0.1--0.9）调节分割敏感度。
\end{itemize}

\textbf{变化检测模式：}
\begin{itemize}
    \item 上传同一地点不同时间的两张图像；
    \item 显示 T1/T2 并列对比、像素差异图、变化概率热力图；
    \item 变化区域半透明覆盖叠加；
    \item 变化面积占比 + 平均/最高置信度统计。
\end{itemize}

模型使用 \lstinline{@st.cache_resource} 实现延迟加载，首次推理后缓存模型实例，后续请求直接复用。

%====================================================================
\section{测试与结果分析（占比20\%）}
%====================================================================

\subsection{功能测试}

\begin{table}[H]
\centering
\caption{功能测试用例}
\label{tab:ft}
\begin{tabular}{clll}
\toprule
编号 & 测试用例 & 预期效果 & 实际结果 \\
\midrule
1 & 上传 PNG 格式遥感图像 & 正常显示并完成分割 & 通过 \\
2 & 上传 JPG 格式遥感图像 & 正常显示并完成分割 & 通过 \\
3 & 上传 TIFF 格式遥感图像 & 正常显示并完成分割 & 通过 \\
4 & 勾选"使用样例"进行分割 & 加载样例图并推理 & 通过 \\
5 & 上传 T1/T2 进行变化检测 & 显示概率图和覆盖图 & 通过 \\
6 & 调节阈值滑块至 0.7 & 分割区域缩小 & 通过 \\
7 & 不传图直接点推理 & 提示"上传图像" & 通过 \\
8 & 上传非图像文件（.txt） & 提示格式不支持 & 通过 \\
\bottomrule
\end{tabular}
\end{table}

\subsection{性能测试}

\begin{table}[H]
\centering
\caption{性能测试数据}
\label{tab:pt}
\begin{tabular}{lcc}
\toprule
测试项 & 语义分割 & 变化检测 \\
\midrule
模型参数量 & 1.97M & 625K \\
GPU 单次推理 & 8.2ms & 12.5ms \\
CPU 单次推理 & 85ms & 110ms \\
GPU 批量推理 (32张) & 120ms & 180ms \\
模型加载时间 & 0.8s & 0.5s \\
训练总时长 (RTX 5060 Ti) & 15min (50 epoch) & 8min (30 epoch) \\
峰值显存占用 & 3.5GB & 1.9GB \\
模型文件大小 & 7.6MB & 2.5MB \\
\bottomrule
\end{tabular}
\end{table}

\subsection{稳定性测试}

\begin{table}[H]
\centering
\caption{稳定性测试}
\label{tab:st}
\begin{tabular}{cllc}
\toprule
编号 & 测试场景 & 系统表现 & 容错措施 \\
\midrule
1 & 传入空白（全黑）图像 & 输出 0\% 建筑占比 & 阈值过滤 \\
2 & 传入全白图像 & 输出 100\% 建筑占比 & 阈值过滤 \\
3 & T1/T2 尺寸不一致 & 自动 Resize 统一 & 预处理阶段处理 \\
4 & 模型文件缺失 & 提示"未训练"并用随机权重 & 异常捕获 \\
5 & 并发多次点击推理 & 队列执行互不干扰 & Streamlit 自动处理 \\
6 & GPU 显存不足 & 自动回退到 CPU & try-except 捕获 \\
\bottomrule
\end{tabular}
\end{table}

%====================================================================
\section{总结与展望（占比10\%）}
%====================================================================

\subsection{项目完成情况}

对照初始目标，项目完成度如下：

\begin{table}[H]
\centering
\caption{目标完成度}
\label{tab:goal}
\begin{tabular}{lccc}
\toprule
目标 & 设定值 & 实际值 & 完成度 \\
\midrule
分割 mIoU & $\geq$85\% & 87.7\% & ✅ 超额完成 \\
检测 F1 & $\geq$60\% & 65.1\% & ✅ 完成 \\
单图推理 & $\leq$50ms & 8.2ms & ✅ 超额完成 \\
Web 界面 & 可实现 & 已完成 & ✅ 完成 \\
\bottomrule
\end{tabular}
\end{table}

\textbf{已解决的实际工程问题：}
\begin{itemize}
    \item 遥感图像跨平台环境配置（Windows/macOS/Linux 一键脚本）；
    \item 轻量级模型在有限参数量下实现高精度分割；
    \item 变化检测中光照/季节噪声与真实变化的区分；
    \item 从数据加载到 Web 部署的完整工程链路。
\end{itemize}

\subsection{问题与解决方案}

\begin{table}[H]
\centering
\caption{关键技术难题与解决方案}
\label{tab:issues}
\begin{tabular}{cll}
\toprule
编号 & 问题 & 解决方案 \\
\midrule
1 & 阿里云镜像断流 & 改用 curl 下载 wheel 再本地 pip 安装 \\
2 & conda run 不支持 Python -c 换行 & 改为写临时 .py 脚本执行 \\
3 & Windows 缺少 Bash 环境 & 提供 setup.bat 双击入口 \\
4 & 数据集来源多样性 & 设计多目录名自动探测机制 \\
5 & 变化检测验证 F1 停滞在 0.35 & 余弦退火降低学习率后提升至 0.65 \\
\bottomrule
\end{tabular}
\end{table}

\subsection{不足与改进方向}

\textbf{现有不足：}
\begin{itemize}
    \item 变化检测 F1 偏低（0.65），误报率较高；
    \item 仅支持二类分割，无法区分多类地物；
    \item 无模型集成机制，单模型结果波动大；
    \item 缺少后处理步骤，存在孤立噪声像素。
\end{itemize}

\textbf{后续优化方向：}

\textbf{短期：}
\begin{itemize}
    \item 增加编码器宽度提升变化检测精度；
    \item 添加 CRF 后处理去除噪点；
    \item 集成 Focal Loss 缓解类别不平衡。
\end{itemize}

\textbf{中期：}
\begin{itemize}
    \item 多任务学习：共享编码器同时输出分割和检测结果；
    \item 半监督预训练（MAE）利用无标注数据；
    \item 添加道路、水体等多类分割。
\end{itemize}

\textbf{长期：}
\begin{itemize}
    \item 集成大模型实现自然语言交互；
    \item ONNX/TensorRT 导出，部署到边缘设备；
    \item 视频遥感序列分析。
\end{itemize}

%====================================================================
\section*{附录}
%====================================================================

\subsection*{参考文献}

\begin{enumerate}
    \item Ronneberger, O., et al. U-Net: Convolutional Networks for Biomedical Image Segmentation. \textit{MICCAI 2015}.
    \item Daudt, R. C., et al. Fully Convolutional Siamese Networks for Change Detection. \textit{ICIP 2018}.
    \item Chen, H., \& Shi, Z. A Spatial-Temporal Attention-Based Method for Remote Sensing Image Change Detection. \textit{Remote Sensing}, 2020.
    \item Loshchilov, I., \& Hutter, F. Decoupled Weight Decay Regularization. \textit{ICLR 2019}.
    \item WHU Building Dataset. \url{http://gpcv.whu.edu.cn/data/}
    \item LEVIR-CD Dataset. \url{https://justchenhao.github.io/LEVIR/}
\end{enumerate}

\subsection*{项目代码仓库}

项目代码已开源：\url{https://github.com/Air-0000/RS_Analysis}

\begin{itemize}
    \item \texttt{REPORT.tex} -- 本报告（LaTeX 源文件）
    \item \texttt{INTRO.md} -- 45KB 技术白皮书
    \item \texttt{setup.sh} / \texttt{setup.bat} -- 一键环境配置
    \item \texttt{app.py} -- Streamlit 交互界面
    \item \texttt{models/} -- 模型定义
    \item \texttt{utils/} -- 数据加载与评估工具
\end{itemize}

\end{document}
